\section{Force: Connections on the Links (P8)}
\label{sec:forces}

A force in Vibe Theory is an instruction carried along a note, the twist applied to a tone as it crosses a link. This is exactly a gauge connection. P8 climbs a ladder: a gauge field on the links (Stage A), coupled to a charged fermion (Stage B), confining (Stage C), and feeling matter and topology (Stage D and beyond).

\subsection{Confinement}

We put an SU(2) gauge field on a periodic lattice, with link variables as unit quaternions and the Wilson action, and measured the static potential through the Creutz ratio (the string tension).

\begin{result}[Non-Abelian confinement]
In 3D SU(2) lattice gauge theory the string tension is positive at every coupling and falls from $1.32$ to $0.40$ as $\beta$ rises, while the average plaquette climbs from $0.10$ to $0.65$. A positive string tension is the Wilson area law, the hallmark of confinement, the defining behaviour of the strong force.
\end{result}

\subsection{The fermion sees gauge topology}

We coupled the chiral overlap fermion to a U(1) gauge background of known topological charge $Q$ and computed its index from the spectral asymmetry of the kernel $H_W$, using a complex-Hermitian eigensolver built for the purpose.

\begin{result}[The lattice index theorem]
The overlap fermion's index equals the gauge topological charge exactly: index $= -Q$, an integer, for every charge tested, with the gauge flux equal to $2\pi Q$ and the kernel Hermiticity error zero. This is the lattice Atiyah-Singer index theorem. The chiral fermion sees gauge topology exactly, at finite lattice spacing.
\end{result}

\subsection{The chiral condensate}

Finally we coupled the overlap fermion to a \emph{dynamical} gauge field, the Schwinger model and its non-Abelian cousin, and measured the chiral condensate through the Banks-Casher near-zero spectral density of the gamma-five-Hermitian overlap.

\begin{result}[A condensate forms from the anomaly]
In a dynamical Abelian gauge field the chiral condensate signal is zero in the free theory and nonzero once the field is on, rising at moderate coupling. The same pattern holds in a dynamical non-Abelian SU(2) field. The chiral fermion feels the fluctuating field and condenses, the anomaly-induced behaviour of the Schwinger model, in both the Abelian and the non-Abelian case.
\end{result}

So the link of Vibe Theory and the gauge connection of physics are the same object. Force is the instruction on the connection, it confines, it carries a topological index that ties it to matter, and it produces a condensate. Two honesty notes belong here. These are small-lattice results by exact diagonalisation: the string tension is the Creutz-ratio estimator rather than the asymptotic value, and the condensate is the Banks-Casher near-zero density (a signal proportional to the condensate) rather than its precise Schwinger value, which would need the dynamical fermion determinant and a continuum extrapolation. And the genuinely open rung, the Weyl-projected non-Abelian chiral gauge theory, has no satisfactory non-perturbative construction anywhere in physics, so we reach the rung below it (a vector-like fermion in the field) rather than the wall itself.
